\section*{Methods}

The general approach is to run classical molecular dynamics simulations for \textit{Z. mobilis} membranes in increasing concentrations of small molecules that might perturb the membrane structure, and would be observed in real hydrolysates.
We also simulate a hypothetical membrane model for \textit{Z. mobilis} that lacks any hopenoids, in order to compare and contrast the impact these small molecules have on membrane structure and dynamics when exposed to multiple stresses.

\subsection*{Membrane Model and System Preparation}
%\todo{This is an unacceptable caption if you want this in a good journal. Figures also should go above where they are cited, as otherwise LaTeX has this nasty habit of moving the damn figures.}
%\todo{Nitin, what are all these abbreviations for the tail lengths? These small molecules aren't named anywhere: Full names are added in table}
\begin{figure}[htbp!]
	\centering
	\includegraphics[width=1.00\linewidth]{Figures/membrane_template1.pdf}
	%\includegraphics[scale=0.50]{Figures/membrane_template1.pdf}

	\caption{(a) Top view of \textit{Z. mobilis} model membrane structure. 
		The membrane is comprised of phosphatidylcholines (PCs) in green, phosphatidylglycerols (PGs) in yellow, phosphoethanolamines (PEs) in blue, cardiolipin (TOCL1) in grey and hopanoids in orange and red color.
		(b) The chemical structures of the lipids, colored according to their head groups. The names for these molecules are given in Table~\ref{tab:lipid_composition}.}

	\label{fig:membrane components}
\end{figure}

%\todo{are these numbers per membrane or per leaflet? Is that clear from the text?}
\begin{table}
	\centering
	\begin{tabular}{p{3.5cm}p{8cm}c}
		\toprule
		Lipid & Name & Numbers \\
		\midrule
		DOPE 18:1/18:1 & 1,2-dioleoyl-\textit{sn}-glycero-3-phosphoethanolamine & 18 \\
		PYPE 16:0/16:1 & 1-palmitoyl-2-arachidonoyl-\textit{sn}-glycero-3-phosphoethanolamine & 4 \\
		YOPE 16:1/18:1 & 1-stearoyl-2-oleoyl-\textit{sn}-glycero-3-phosphoethanolamine & 3 \\
		POPE 16:0/18:1 & 1-palmitoyl-2-oleoyl-\textit{sn}-glycero-3-phosphoethanolamine & 3 \\
		DYPE 16:1/16:1 & 1,2-di-(9\textit{Z}-hexadecenoyl)-\textit{sn}-glycero-3-phosphoethanolamine & 1 \\
		\midrule
		DOPG 18:1/18:1 & 1,2-dioleoyl-\textit{sn}-glycero-3-phospho-(1$'$-rac-glycerol) & 8 \\
		POPG 16:0/18:1 & 1-palmitoyl-2-oleoyl-\textit{sn}-glycero-3-phospho-(1$'$-rac-glycerol) & 4 \\
		YPPG 16:1/16:0 & 1-palmitoleoyl-2-palmitoyl-\textit{sn}-glycero-3-phospho-(1$'$-rac-glycerol) & 1 \\
		\midrule
		DOPC 18:1/18:1 & 1,2-dioleoyl-\textit{sn}-glycero-3-phosphocholine & 4 \\
		POPC 16:0/18:1 & 1-palmitoyl-2-oleoyl-\textit{sn}-glycero-3-phosphocholine & 1 \\
		PYPC 16:0/16:1 & 1-palmitoyl-2-palmitoleoyl-\textit{sn}-glycero-3-phosphocholine & 1 \\
		\midrule
		TOCLI 18:1/18:1/18:1/18:1 & 1$'$,3$'$-bis(1,2-dioleoyl-\textit{sn}-glycero-3-phospho)-\textit{sn}-glycerol & 2 \\
		\midrule
		& Bacteriohopanetetrol cyclitol ether & 25 \\
		& Bacteriohopanetetrol glucosamine  & 25 \\
		\midrule
		Total & & 100 \\
		\bottomrule
	\end{tabular}
	\caption{Membrane phospholipid and hopanoid composition in wild type \emph{Z. mobilis} adapted from available lipidomics\cite{RiveraVazquez2024} to specify the number of lipids in each leaflet of the lipid bilayer.}
	\label{tab:lipid_composition}
\end{table}


The \textit{Z. mobilis} membrane model was constructed from the phospholipids phosphatidylethanolamine (PE), phosphatidylglycerol (PG), phosphatidylcholine (PC), and cardiolipin (CL), alongside hopanoids (bacteriohopanetetrol cyclitol ether and bacteriohopanetetrol glucosamine), as illustrated in Figure~\ref{fig:membrane components}.
The composition of the model was based on recent lipidomics data\cite{RiveraVazquez2024}, condensing the multitude of lipids identified by lipidomics into a representative lipid mixture as detailed in Table~\ref{tab:lipid_composition}.
Each membrane leaflet contained 100 molecules, consisting of either a mixture of phospholipids and hopanoids or phospholipids alone.
For the hypothetical hopanoid-deficient membrane system, the membrane was generated using the CHARMM-GUI Membrane Builder\cite{Lee2019b}. Each leaflet contained 100 phospholipid molecules (combining PE, PG, PC, and CL), utilizing double the quantities listed in Table~\ref{tab:lipid_composition} to compensate for the absence of hopanoids.

The hopanoid-rich membrane system was constructed through a similar multi-step process. 
Initially, the system was built using the CHARMM-GUI Membrane Builder\cite{Lee2019b}. Because hopanoids are not currently present in the CHARMM-GUI lipid library, cholesterol and ergosterol were used as placeholders for the 50 hopanoid molecules required per leaflet. 
This 1:1 ratio represents the hopanoid enrichment found in \emph{Zymomonas} strains actively producing ethanol\cite{Hermans1991}, as hopanoids are essential to microbial fitness under these conditions\cite{Brenac2019}.
Following initial membrane assembly, each cholesterol or ergosterol molecule was replaced with the specific hopanoids found in \textit{Z. mobilis} membranes\cite{Brenac2019}. 
Finally, as this replacement process can result in ring penetrations, LongBondEliminator was utilized to resolve these artifacts prior to extended simulation\cite{Sarkar2023}.




%\begin{comment}
%\begin{table}
%	\centering
%	\begin{tabular}{ll}
%		\toprule
%		\textbf{Class} & \textbf{Compounds} \\
%		\midrule
%		Carboxylic Acids &  Formic Acid, Acetic Acid, Propanoic Acid \\
%		Alcohols & Ethanol,  Propanol, Isobutanol \\
%		Aldehydes & Acetaldehyde \\
%		Ketones & Acetone \\
%		Furan Compounds & HMF (5-Hydroxymethylfurfural), Furfural \\
%		\bottomrule
%	\end{tabular}
%	\caption{Classification and list of small molecules added to the membrane at different concentrations.}
%	\label{tab:compound_classes}
%\end{table}
%\end{comment}


To study the interactions of small molecules (Figure~\ref{fig:molecules_images}) with the \textit{Z. mobilis} membrane, simulations were performed at varying concentrations of the small molecules (0, 0.5, 1.0, 1.5, 2.0, and 2.5 mol\%).
Since the carboxylic acids are less hydrophilic than their charged conjugate base, and that the acids are the primary form that crosses a typical lipid bilayer\cite{Lee2016,Vermaas2019a,Raza2023}, we only simulate the acid in this study.
All small molecules were initially placed in the water phase, as shown in Figure~\ref{fig:system_setup}.
The system was solvated with TIP3P water molecules and was neutralized by adding counterions (\ce{Na+} or \ce{Cl-}).


%\todo{Hey Nitin, where is the description of how you made the systems without hopenoids?: Now we have it}

\subsection*{Molecular Dynamics Simulations}
%We performed classical MD simulations with the classical CHARMM36 for lipids\cite{Klauda2010} and CGenFF for the small molecules\cite{Vanommeslaeghe2010}.
%For each concentration built above, three independent replicas were simulated for adaquate statistical sampling.
All MD simulations were performed using the NAMD 3.0.1 simulation engine \cite{Phillips2020}. 
The lipid parameters were derived from CHARMM36~\cite{Yu2020}, while the CGenFF~\cite{Vanommeslaeghe2010} was employed for the small molecules and hopanoids (Bacteriohopanetetrol cyclitol ether and Bacteriohopanetetrol glucosamine).
This classical molecular simulation captures partitioning, but not reactivity for molecules like aldehydes present in our molecule set~\cite{Farah2012}.
As is standard for CHARMM36\cite{Yu2020}, we used the TIP3P water model\cite{Jorgensen1983}, the representative initial system setup is shown in Figure~\ref{fig:system_setup}.
Periodic boundary conditions were applied in all directions to avoid edge effects. 
The system comprised 100 lipids per leaflet and an initial simulation box of approximately $89 \times 93 \times 100~\text{\AA}^3$. 
The lateral dimensions were chosen to provide sufficient in-plane extent to minimize finite-size effects in heterogeneous membranes, while remaining consistent with established all-atom simulations of multi-component lipid bilayers\cite{Gupta2016,Choe2024,Rice2025}.

%JVV, \, gives you a thin space.
The systems were first energy-minimized using the steepest descent algorithm to remove any steric clashes.
The simulations were performed in a constant pressure and temperature (NPT) ensemble. 
The temperature was maintained at 300\,K using the Langevin thermostat\cite{Dunweg2003}, and the pressure was controlled at 1 bar using the Langevin barostat\cite{Feller1995}.
The barostat had a piston period of 200 fs and a decay time of 100\,fs.
Group pressure coupling was used in combination with a flexible cell and a constant cell ratio, as is typical for membrane simulations in NAMD.
The flexible cell setting permits the simulation box to adjust both in volume and shape to accommodate anisotropic pressure fluctuations, while the constant ratio between the x- and y- dimensions keeps the membrane aspect ratio constant.
A 2\,fs time step was applied throughout the simulation, enabled by constraining bond lengths to hydrogen atoms using SETTLE\cite{Miyamoto1992}.
Long-range electrostatic interactions were calculated using the Particle Mesh Ewald (PME) method \cite{Essmann1995} with a grid spacing of 1\,\Angs.
Van der Waals interactions were truncated at 12\,\Angs with a switching function applied at 10\,\Angs to maintain continuous forces and energies.
All the simulations were performed for three independent runs, each run to 1000~ns.
All thermodynamic analysis was carried out for the last 800\,ns considering the first 200\,ns as the equillibration time, and statistics were averaged over the three independent runs.

\subsection*{Membrane Property Analysis}
We measure multiple membrane properties from the simulations described above.
All properties were calculated for each of the three replicas, and the results were averaged over all replicas.
Error estimates were determined as the standard error of the mean across the average of individual simulation replicates.
All calculations were performed using in-house Python scripts, using what numpy\cite{vanderWalt2011}, scikit\cite{JMLR:v12:pedregosa11a}, and matplotlib\cite{Hunter2007}.
The following properties were analyzed to characterize the interactions of small molecules with the \textit{Z. mobilis} membrane:

\subsubsection*{Area per Lipid (APL)}
The area per lipid (APL) was determined from the periodic box size of the simulation system, averaged over the equilibrated portion of the trajectory and divided by the fixed number of lipids in a leaflet.
While we could have determined the area per lipid for individual lipids, we focus here on the overall area per lipid to track changes in membrane structure when the small molecule concentration rises.
% as the average surface area of the membrane divided by the number of lipids in one leaflet. The membrane surface area was determined from the dimensions of the simulation box in the xy-plane.

\subsubsection*{Membrane Thickness}

Similarly, we measure the membrane thickness for the entire membrane by measuring the distance along the membrane normal (z) axis between the average position for phosphorus atoms in the upper and lower leaflets.
%The thickness was calculated along the membrane normal-axis. It was measured as the distance between the Phosphate atoms of the lipids in the upper and lower leaflets. 


\subsubsection*{Lipid Order Parameter ($S_{CH}$)}
To facilitate comparisons with potential future NMR experiments, we determine the lipid order parameter, averaged across all lipids all along every acyl chain to measure the overall order as a function of small molecule concentration.
The lipid order parameter, $S_{CH}$, is measured in NMR via carbon-deuterium bond orientations, but we are measuring based on the equivalent C-H bond vector.
The angle of this bond compared with the membrane normal axis, $\theta$, is used to quantify lipid order:

\begin{equation}
	\label{eq:order}
S_{CH} = \frac{1}{2} \langle 3 \cos^2 \theta - 1 \rangle
\end{equation}

\subsubsection*{Lateral Diffusion Coefficient}
Measuring overall lipid dynamics is done here by measuring the lateral diffusion of individual lipids within the larger molecular system.
The lipid lateral diffusion coefficient ($D_{xy}$) was calculated from the mean square displacement (MSD) of lipids in the xy-plane using the relation:
\begin{equation}
	\label{eq:diffusion}
D_{xy} = \frac{\langle MSD_{xy} \rangle}{4}
\end{equation}
%\todo{Did you fit a line to compute a slope? Or were you actually using lag times? If so, what was the lag time? This is insufficiently clear}

To compute $D_{xy}$, the molecular dynamics trajectory was divided into consecutive 10~ns segments. For each segment, the MSD of the lipid phosphorus atoms was calculated relative to the first frame of that segment.
A linear regression was then performed on the MSD versus time data for each segment, and the fitted slope was used to determine $D_{xy}$ according to Equation~\ref{eq:diffusion}. 
The final diffusion coefficient was reported as the mean of all segment-based values, with the standard error of the mean calculated to represent variability across segments.

\subsubsection*{Probability Density and Free Energy Profile for Small Molecules Across the Membrane}
From our extensive sampling, it is possible to measure the probability distribution for the various small molecules across the membrane.
This is done by determining the center of mass for each molecule with respect to the membrane normal, and then histogramming the resulting positions to determine a probability distribution.
The probability distribution can then be converted into a free energy profile by the well worn relationship:
\begin{equation}
	\label{eq:boltzmann}
	\Delta G = -RT \ln \frac{p}{p_0} = G_p - G_{p_0}
\end{equation}
In Eq.~\ref{eq:boltzmann}, p is the probability within a specific bin, while $p_0$ is the probability in a reference bin within the histogram.
We think it is most helpful if the small molecule in solution is considered to be the zero point for free energy, which we accomplish by selecting the probability in solution as $p_0$.
%The probability density of ethanol and acetic acid molecules along the membrane normal (z-axis) was calculated to determine their preferred localization within the membrane. 


\subsubsection{Flux-based estimation of permeability}

The permeability values were calculated
directly from the statistics of complete leaflet-to-leaflet crossings observed
in atomistic MD trajectories.
For a given solute we recorded the total number
of full translocation events~$N$ in three independent $1000\;\text{ns}$
replicas of a hydrated \emph{Z. mobilis} model membrane.  
A permeability coefficient was then
computed following the analysis framework from \citet{Venable2019}.
%
\begin{equation}
	Pm = \frac{N}{2 A \, \Delta t \, C}
	\label{eq:permeability}
\end{equation}
%
where the factor~$2$ accounts for the two bilayer faces, $\Delta t$ is the
aggregate simulation time and $C$ is the instantaneous aqueous concentration
extracted from the bulk region and \emph{A} is the area of the lipid bilayer.